% Part: first-order-logic
% Chapter: axiomatic-deduction
% Section: deduction-theorem

% verification of properties of provability needed for maximally
% consistent sets in the completeness chapter.

\documentclass[../../../include/open-logic-section]{subfiles}

\begin{document}

\iftag{FOL}
      {\olfileid{fol}{axd}{ded}}
      {\olfileid{pl}{axd}{ded}}

\olsection{مبرهنة الاستنباط}

كما رأينا، فإن إعطاء !!{derivation}s في نسق بديهي أمر شاق، وقد يصعب
العثور على !!{derivation}s. وبدلًا من أن نكتب فعلًا قوائم طويلة من
ال!!{formula}s، يسهل بوجه عام أن نبرهن على وجود مثل هذه
ال!!{derivation}s بالاستعانة ببعض النتائج البسيطة. وقد أثبتنا بالفعل
ثلاثًا من هذه النتائج: تفيد \olref[ptn]{prop:reflexivity} بأنه يمكننا
دائمًا الجزم بأن $\Gamma \Proves !A$ متى علمنا أن $!A \in
\Gamma$. وتفيد \olref[ptn]{prop:monotonicity} بأنه إذا كان $\Gamma \Proves !A$
كان أيضًا $\Gamma \cup \{!B\} \Proves !A$. ويترتب من
\olref[ptn]{prop:transitivity} أنه إذا كان $\Gamma \Proves !A$ و
$!A \Proves !B$، فإن $\Gamma \Proves !B$. وهذه نتيجة بسيطة أخرى،
وهي نسخة «فوقية» من قاعدة إثبات المقدَّم:

\begin{prop}
  \ollabel{prop:mp} إذا كان $\Gamma \Proves !A$ و$\Gamma \Proves !A \lif
  !B$، فإن $\Gamma \Proves !B$.
\end{prop}

\begin{proof}
  لدينا $\{!A, !A \lif !B\} \Proves !B$:
  \begin{derivation}
    1. & $!A$ & \Hyp\\
    2. & $!A \lif !B$ & \Hyp\\
    3. & $!B$ & 1، 2، إثبات المقدَّم
  \end{derivation}
  وبحسب \olref[ptn]{prop:transitivity}، نحصل على $\Gamma \Proves !B$.
\end{proof}

أما أهم نتيجة سنستعملها في هذا السياق فهي مبرهنة الاستنباط:

\begin{thm}[مبرهنة الاستنباط]
\ollabel{thm:deduction-thm} $\Gamma \cup \{!A\} \Proves !B$ إذا وفقط
إذا كان $\Gamma \Proves !A \lif !B$.
\end{thm}

\begin{proof}
اتجاه «إذا» مباشر. فإذا كان $\Gamma \Proves !A \lif !B$، كان أيضًا
$\Gamma \cup \{!A\} \Proves !A \lif !B$ بحسب
\olref[ptn]{prop:monotonicity}. وكذلك $\Gamma \cup \{!A\} \Proves !A$
بحسب \olref[ptn]{prop:reflexivity}. ومن ثم تعطينا \olref{prop:mp} أن
$\Gamma \cup \{!A\} \Proves !B$.

أما اتجاه «فقط إذا» فنبرهنه بالاستقراء في طول !!{derivation} الصيغة
$!B$ من $\Gamma \cup \{!A\}$.

في أساس الاستقراء نبرهن الدعوى لكل !!{derivation} طوله~$1$.
يتألف !!^a{derivation} للصيغة~$!B$ من $\Gamma \cup \{!A\}$ وطوله~$1$
من $!B$ وحدها؛ وإذا كان صحيحًا، فإما أن تكون
$!B \in \Gamma \cup \{!A\}$ أو أن تكون بديهية. وإذا كان $!B \in \Gamma$
أو كان بديهية، فإن $\Gamma \Proves !B$. ولدينا أيضًا $\Gamma
\Proves !B \lif (!A \lif !B)$ بحسب \olref[prp]{ax:lif1}، ومن
\olref{prop:mp} نحصل على $\Gamma \Proves !A \lif !B$. وإذا كان
$!B \in \{ !A\}$، فإن $\Gamma \Proves !A \lif !B$، لأن ال!!{formula}
الأخيرة~$!A \lif !B$ هي نفسها $!A \lif !A$، وقد قدّمنا !!{derive}
لهذه الصيغة في \olref[pro]{ex:identity}.

وفي خطوة الاستقراء، افترض أن !!a{derivation} للصيغة~$!B$ من
$\Gamma \cup \{!A\}$ ينتهي بخطوة~$!B$ تبررها قاعدة إثبات المقدَّم.
(فإن لم تبررها قاعدة إثبات المقدَّم، كان $!B \in \Gamma$ أو $!B
\ident !A$ أو كانت $!B$ بديهية، وينطبق عندئذ التعليل نفسه المستعمل في
أساس الاستقراء.) إذن من الخطوات السابقة في ال!!{derivation} الصيغتان
$!C \lif !B$ و$!C$، لبعض ال!!{formula}~$!C$؛ أي إن
$\Gamma \cup \{!A\} \Proves !C \lif !B$ و$\Gamma \cup \{!A\}
\Proves !C$. وال!!{derivation}s المقابلة أقصر، فينطبق عليها فرض
الاستقراء. ومن ثم نحصل على كلتا العبارتين:
\begin{align*}
  & \Gamma \Proves !A \lif (!C \lif !B); \\
  & \Gamma \Proves !A \lif !C.
\end{align*}
ولدينا أيضًا
\[
\Gamma \Proves (!A \lif (!C \lif !B)) \lif
((!A\lif !C)  \lif (!A \lif !B)),
\]
بحسب \olref[prp]{ax:lif2}، ويعطينا تطبيقان لـ\olref{prop:mp} النتيجة
$\Gamma \Proves !A \lif !B$، كما هو مطلوب.
\end{proof}

لاحظ أن \olref[prp]{ax:lif1} و\olref[prp]{ax:lif2} قد اختيرتا على وجه
التحديد بحيث تصح مبرهنة الاستنباط.

وفيما يلي بعض الحقائق المفيدة عن !!{derivability}، ونترك برهانها تمارين.

\begin{prop}
  \ollabel{prop:derivfacts}
\begin{enumerate}
\item $\Proves (!A \lif !B) \lif ((!B \lif !C)
  \lif (!A \lif !C))$؛ \ollabel{derivfacts:a}
\item إذا كان $\Gamma \cup \{ \lnot !A\}
  \Proves \lnot !B$، فإن $\Gamma \cup \{ !B\} \Proves
  !A$ (المقابلة)؛ \ollabel{derivfacts:b}
\item  $\{ !A, \lnot!A\} \Proves
    !B$ (من الباطل يلزم أي شيء، الانفجار)؛ \ollabel{derivfacts:c}
\item  $\{ \lnot\lnot!A\} \Proves
  !A$ (حذف النفي المزدوج)؛\ollabel{derivfacts:d}
\item إذا كان $\Gamma \Proves \lnot\lnot!A$، فإن $\Gamma \Proves
  !A$؛\ollabel{derivfacts:e}
\end{enumerate}
\end{prop}

\tagprob{FOL}
\begin{prob}
  أثبت \olref[fol][axd][ded]{prop:derivfacts}.
\end{prob}
\tagendprob

\tagprob{notFOL}
\begin{prob}
  أثبت \olref[pl][axd][ded]{prop:derivfacts}.
\end{prob}
\tagendprob

\end{document}
